\paragraph{The lattice estimate for the CM approximations.}

The lower bound for the Chow estimate uses an integral section evaluated near a CM frame. Hermitian transference controls how closely the required directions can be approximated.

\begin{lemma}[Hermitian lattice transference]
\label{lem:cm-lattice-transference}
Let \(E\) be the CM field, \(r_0\ge1\) an integer, and \(L\subset E^{r_0}\) a full \(\mathcal O_E\)-lattice
self-dual for a Hermitian form \(H\). Take one embedding from
each conjugate pair and write \(f=[E:\mathbb Q]/2\).
At these embeddings let \(Q_\tau\) take \(H_\tau\) to its
diagonal signature matrix \(J_\tau\). Write
\(q_\tau(z)=\|Q_\tau\tau(z)\|^2\) for the corresponding
positive definite majorants. Suppose

\[
  \max_\tau q_\tau(z)\ge\lambda^2>0
\]
 for every nonzero
\(z\in L\). Put \(d_E=|\operatorname{disc}E|\),
\(t_*=fr_0\), and use the real Euclidean space
\(\prod_\tau\mathbb C^{r_0}\) of dimension \(2t_*\).
Then the image lattice under \(z\mapsto(Q_\tau\tau(z))_\tau\)
has covering radius at most \(4t_*^3d_E/\lambda\). There is
an \(E\)-basis of \(E^{r_0}\) in \(L\) with all Euclidean
image norms at most the same bound.

\end{lemma}

\begin{proof}
\proofstep{Comparing the rational and Hermitian duals.}
Indeed let \(L^\dagger\) be the rational dual of \(L\) for
the pairing \(\operatorname{Tr}_{E/\mathbb Q}H\). For
\(z\in L^\dagger\), pairing with all \(\mathcal O_E\)-multiples
of a vector of \(L\) puts each \(H(z,y)\) in the inverse different.
Thus \(d_E z\in L\) by Hermitian self-duality. Combining a CM type
with its conjugate shows that the image of \(L^\dagger\) under
\(z\mapsto(2J_\tau Q_\tau\tau(z))_\tau\) is the Euclidean dual.

In fact the Euclidean pairing of these two images is
\(2\sum_\tau\operatorname{Re}\tau H(z,y)=\operatorname{Tr}H(z,y)\).
The factor two identifies a real dual basis with the image of the
rational trace dual basis. For \(0\ne z\in L^\dagger\), the element
\(d_Ez\) is a nonzero vector of \(L\), whence
\[
 \Big(\sum_\tau\|2J_\tau Q_\tau\tau(z)\|^2\Big)^{1/2}
 \ \ge\ \frac{2}{d_E}\max_\tau\sqrt{q_\tau(d_Ez)}
 \ \ge\ \frac{2\lambda}{d_E}.
\] Euclidean transference gives \(2t_*\)
independent primal vectors of norm at most
\((2t_*)^2/(2\lambda/d_E)=2t_*^2 d_E/\lambda\).
Their integer span is a sublattice.

For an arbitrary point write \(x=\sum_{j=1}^{2t_*}a_j v_j\) in
this real basis, and choose \(n_j\in\mathbb Z\) with
\(|a_j-n_j|\le1/2\). Then
\[
 \Big\|x-\sum_j n_jv_j\Big\|
 \le\tfrac12\sum_j\|v_j\|
 \le \tfrac12(2t_*)\,2t_*^2d_E/\lambda=2t_*^3 d_E/\lambda .
\]
The rounded point belongs to the original lattice. The same family spans \(E^{r_0}\)
over \(\mathbb Q\) after pulling back, so an \(E\)-basis can be
selected from it. Both conclusions fit the stated bound.
This argument is used in the full dimension and on the
complementary lattice when it has the same self-duality and
majorant hypotheses. The quantitative choice of the frame
uses the resulting covering radius at that frame's cusp times.
\cite{banaszczyk}\label{cor:slice-cm-transference-1}
\end{proof}
